% !TEX root = ../main.tex
%% Devoir surveillé n° 1 — chapitres 1 à 3
%% Résultats vérifiés par verifications/eval_ds.py.

\sujetEval{DS1}{MPSI · Devoir surveillé}{Devoir surveillé \numero1}
  {Durée : \textbf{4 heures}\\Calculatrice interdite\\Chapitres 1 à 3}{4}
  {Exercice 1/Ensembles et applications/5 pts,Exercice 2/Sommes et récurrences/5 pts,%
   Exercice 3/Trigonométrie/4 pts,Exercice 4/Nombres complexes/6 pts}
\consigne{Les quatre exercices sont indépendants. La qualité de la rédaction
et la précision des raisonnements entrent pour une part importante dans
l'évaluation ; les résultats seront encadrés.}

\begin{exoEval}{5 points}{Ensembles et applications}
On considère $f:\R\to\R$, $x\mapsto x^{2}-2x$.
\begin{questions}
  \item Déterminer $f\bigl(\intFF{-1}{2}\bigr)$ et
    $f^{-1}\bigl(\intFF{0}{3}\bigr)$.
  \item L'application $f$ est-elle injective ? surjective ? Montrer que $f$
    induit une bijection $g$ de $\intFO{1}{+\infty}$ sur
    $\intFO{-1}{+\infty}$ et expliciter $g^{-1}$.
  \item Soient $E,F$ deux ensembles, $h:E\to F$ et $A,B\subseteq E$. Montrer
    que $h(A\cap B)\subseteq h(A)\cap h(B)$, et que l'égalité a lieu pour
    toutes parties $A,B$ lorsque $h$ est injective. Donner, avec la fonction
    $f$, un exemple où l'inclusion est stricte.
\end{questions}
\end{exoEval}

\begin{exoEval}{5 points}{Sommes et récurrences}
Soit $n\in\N^{*}$.
\begin{questions}
  \item Calculer $\displaystyle\sum_{k=1}^{n}k\cdot k!$ \emph{(écrire
    $k=(k+1)-1$)}.
  \item Montrer que $\dfrac{1}{k+1}\dbinom nk=\dfrac{1}{n+1}\dbinom{n+1}{k+1}$,
    et en déduire $\displaystyle\sum_{k=0}^{n}\frac{1}{k+1}\binom nk$.
  \item Calculer la somme double $\displaystyle\sum_{1\leqslant i\leqslant j\leqslant n}i$.
  \item Montrer par récurrence que pour tout réel $x\geqslant-1$ :
    $(1+x)^{n}\geqslant1+nx$ (inégalité de Bernoulli).
\end{questions}
\end{exoEval}

\begin{exoEval}{4 points}{Trigonométrie}
\begin{questions}
  \item Résoudre dans $\intFO{0}{2\pi}$ l'équation
    $2\cos^{2}x+3\sin x=3$.
  \item Calculer la valeur exacte de $\cos\frac\pi8$.
  \item Établir $\cos3\theta=4\cos^{3}\theta-3\cos\theta$. En déduire que
    $\cos\frac\pi9$, $\cos\frac{5\pi}9$ et $\cos\frac{7\pi}9$ sont les
    trois racines de $8t^{3}-6t-1=0$, puis calculer leur somme et leur
    produit.
\end{questions}
\end{exoEval}

\begin{exoEval}{6 points}{Nombres complexes}
\begin{questions}
  \item Résoudre dans $\Cx$ l'équation $z^{2}-(3+\mathrm{i})z+4+3\mathrm{i}=0$.
  \item Pour $z\neq-1$, on pose $Z=\dfrac{z-\mathrm{i}}{z+1}$. En écrivant
    $z=x+\mathrm{i}y$, déterminer et représenter l'ensemble des points
    d'affixe $z$ tels que $Z$ soit réel, puis tels que $Z$ soit imaginaire
    pur.
  \item Montrer que pour tout $z\in\mathbb{U}\setminus\ens{1}$,
    $\dfrac{1+z}{1-z}$ est imaginaire pur, et l'exprimer à l'aide de
    l'argument de $z$.
  \item Soit $n\geqslant2$. Résoudre dans $\Cx$ l'équation
    $(z+1)^{n}=(z-1)^{n}$. Combien a-t-elle de solutions ? Où sont-elles
    situées ?
\end{questions}
\end{exoEval}
